% !PRINTED 521
\weberpage{521}
%<<CONT
und das weder mit $A$ noch mit $B$ ein Element gemein hat. So
f\"ahrt man fort, bis alle Elemente ersch\"opft sind, und kommt zu
folgendem Satze:

\gesperrt{2. Eine transitive Gruppe hat nur wenn sie imprimitiv
 ist, einen von der identischen Gruppe
verschiedenen intransitiven Normaltheiler. Die
Systeme der Intransitivit\"at des Normaltheilers
sind Systeme der Imprimivit\"at der gegebenen
Gruppe.} % [sic] printed \textquotedbl{}Imprimivit\"at\textquotedbl{} (expected \textquotedbl{}Imprimitivit\"at\textquotedbl{})

In Bezug auf die Gleichungen k\"onnen wir mit R\"ucksicht auf
1. den Satz 2. auch so aussprechen:

\gesperrt{3. Eine irreducible Gleichung $f(x) = 0$ zerf\"allt,
wenn sie durch Adjunction aller Wurzeln einer
Resolvente oder einer Wurzel einer Normalgleichung
 $j^{\mathrm{ten}}$ Grades reducibel wird, in mehrere
irreducible Factoren von gleichem Grade und
von gleicher Gruppe. Die Anzahl dieser Factoren
ist ein Theiler des Grades $j$ und ist gleich $j$,
wenn $j$ eine Primzahl ist.}

Wenn also eine irreducible primitive Gleichung durch
Adjunction der Wurzeln einer Resolvente reducirt wird, so zerf\"allt
sie in lineare Factoren, d.\,h. sie ist vollst\"andig gel\"ost.
Dieser Fall tritt immer ein, wenn der Grad der Gleichung
$f(x) = 0$ eine Primzahl ist, weil eine solche Gleichung nicht
imprimitiv sein kann.

\begin{center}
\rule{3cm}{0.4pt}
\end{center}
